\section{The group and an action of small cost}
\label{actions:section}

We construct the group and two free probability-measure-preserving actions
that will be compared.  The purpose of this section is to prove the upper
bound for the action with an additional finite height coordinate.
All actions are right actions: $(zg)h=z(gh)$.

\subsection{The amalgam}

Let $A=F(a,b_1,\ldots,b_{99})$ be the free group on the
displayed generators.  Define words in $A$ by
\[
 u_0=1,\qquad u_i=u_{i-1}ab_i\quad(1\leq i\leq99),
 \qquad w=u_{99}a.
\]
Let
\begin{equation}
 \label{actions:group}
 J=\langle b_1,\ldots,b_{99},w\rangle\leq A,
 \qquad B=J\times\langle t\rangle,
 \qquad \Gamma=A*_J B,
\end{equation}
where $\langle t\rangle$ is infinite cyclic and the embedding of $J$ in
$B$ is $j\mapsto(j,1)$.

\begin{lemma}
\label{actions:normal-form}
The factors $A$ and $B$ embed in $\Gamma$ with intersection $J$.
If $g_1,\ldots,g_n$ lie alternately in $A\setminus J$ and
$B\setminus J$, then $g_1\cdots g_n\notin J$ for $n\geq1$, and
$g_1\cdots g_n$ belongs to neither factor for $n\geq2$.
The group $J$ is free on $b_1,\ldots,b_{99},w$; in particular it is
infinite.  The group $\Gamma$ is countably infinite and generated by
$a,b_1,\ldots,b_{99},t$.
\end{lemma}

\begin{proof}
Choose, in each factor, representatives for the cosets $gJ$, with $1$
representing $J$.  Consider formal strings
\[
 r_1\cdots r_n j,
 \qquad j\in J,
\]
where the $r_i$ are nonidentity representatives, alternately from the
two factors; the case $n=0$ is allowed.  Define right multiplication by
$g$ in one factor as follows.  If the last representative is from that
factor, combine it with the final $j$ and $g$, and write their product
uniquely as $rj'$ in that factor.  Otherwise write $jg=rj'$ and append
$r$.  Omit $r$ when $r=1$.

This gives a right action of each factor on the strings.  To check the
action law for two successive multipliers in the same factor, retain
the prefix ending in the other factor and simply multiply the remaining
suffix in the chosen factor.  This description continues to apply when
the suffix representative becomes $1$.  The two actions agree for
multipliers in $J$, which only multiply the final $J$-entry.  They
therefore define an action of the amalgam on the strings.

Acting on the identity string by a factor element recovers its coset
decomposition, so the factors embed.  Their possible strings intersect
exactly in the strings with no representatives, proving $A\cap B=J$.
Successively multiplying by an alternating sequence outside $J$
adds one nonidentity representative at every step: for $j\in J$ and
$g\notin J$, one has $jg\notin J$.  The resulting string has length
$n$, whereas elements of $J$ have length zero and factor elements have
length at most one.  This proves the stated normal-form assertions.

For $J$, expand any reduced nonempty word in the abstract generators
$b_1,\ldots,b_{99},w$ in the free basis of $A$.
The word $w=ab_1ab_2\cdots ab_{99}a$ starts and ends with $a$;
$w^{-1}$ starts and ends with $a^{-1}$.  No boundary between a
$b_i^{\pm1}$ and $w^{\pm1}$ cancels.  Two successive $w$-symbols
either have equal signs, in which case their boundary does not cancel,
or would be inverse symbols, which reducedness excludes.  Boundaries
between $b$-symbols do not cancel either.  Thus the expansion is a
nonempty reduced word in $A$, proving that the displayed generators of
$J$ are a free basis.

Finally, the displayed generators of $\Gamma$ generate both factors;
the embedded copy of $A$ makes $\Gamma$ infinite, and finite generation
makes it countable.  The embedded copy of $B$ also shows that $t$ has
infinite order in $\Gamma$.
\end{proof}

\subsection{Probability spaces and the two actions}

We use the definitions of graphing, cost, and p.m.p. action from the
introduction. A standard Borel space is a measurable space isomorphic
to a Borel subset of a complete separable metric space. For a subgroup
$L\leq\Gamma$, write $\mathcal R_L$ for its orbit relation on the
space being considered.

For an essentially free action of a countable group, the set of points with trivial stabilizer is a Borel invariant
conull set.  Restricting
to this set makes the action pointwise free and preserves the measure
and orbit-relation cost.  More generally, a Borel conull set can be
made invariant by intersecting its countably many translates.
To see that an invariant conull restriction preserves cost, restrict
any orbit graphing to that set; invariance keeps its edges inside
the set, and domain measures do not change.  Conversely, a graphing
on the restricted set generates the original relation on a conull
set and has the same cost when viewed on the original space.

Let
\[
 X_0=[0,1]^\Gamma,\qquad
 \mu_0=\operatorname{Leb}^{\Gamma},\qquad
 (xg)(h)=x(gh)\quad(g,h\in\Gamma).
\]
This is a standard Borel probability space, since it is a countable
product of compact metrizable spaces, with its product measure.
Coordinate permutations are Borel automorphisms preserving $\mu_0$,
and the displayed formula satisfies the right-action law.  For
$g\ne1$, the equality $xg=x$ implies $x(g)=x(1)$.  These are two
independent uniform coordinates, so this equality has probability zero.
Consequently
\[
 X=\{x\in X_0:xg\ne x\text{ for every }g\ne1\}
\]
is an invariant Borel conull set.  Write $\mu$ for the restricted
probability measure.  We henceforth use the free Bernoulli action on
$(X,\mu)$; it has the same orbit-relation cost as the original action
on $(X_0,\mu_0)$.

\begin{lemma}
\label{actions:height}
There is a homomorphism $\chi:\Gamma\to\mathbb Z$ with
\[
 \chi(a)=1,\qquad \chi(b_i)=\chi(t)=0.
\]
It satisfies $\chi(u_i)=i$ and $\chi(w)=100$.
For every positive integer $M$, the action on
\[
 Y_M=X\times\mathbb Z/M\mathbb Z,
 \qquad \nu_M=\mu\times\text{uniform measure},
\]
defined by
\begin{equation}
 \label{actions:product-action}
 (x,r)g=(xg,r+\chi(g)\bmod M)
\end{equation}
is free and p.m.p. on a standard Borel probability space.
\end{lemma}

\begin{proof}
The exponent sum of $a$ defines a homomorphism $\chi_A:A\to\mathbb Z$.
Define $\chi_B(jt^n)=\chi_A(j)$ on $B=J\times\langle t\rangle$.
This is a homomorphism and agrees with $\chi_A$ on $J$, so the
amalgam's universal property gives $\chi$.  The formulas for $u_i$
and $w$ follow by their definitions.  In particular, the restriction
of $\chi$ to $J$ need not be zero: it sends $w$ to $100$.
The homomorphism law verifies the right-action law in
\eqref{actions:product-action}.  Both coordinate actions preserve
their measures.  Freeness follows from freeness on the first
coordinate, and the finite product remains standard Borel.
\end{proof}

\subsection{A graphing on the height extension}

The direct-product factor $B$ can be generated at arbitrarily small
cost above one by retaining all $t$-edges and using the generators of
$J$ on a small set meeting every $t$-orbit.  This is the product
construction in Gaboriau's Proposition~VI.23
\cite[Proposition~VI.23]{gaboriau-2000}.  For our Bernoulli action the
small set has the following explicit description.

\begin{lemma}
\label{actions:small-marker}
For every $0<p<1$ and every positive integer $M$, there is a Borel
set $C_p\subseteq Y_M$ of measure $p$ meeting every
$\langle t\rangle$-orbit.  The graphing consisting of $t$ on all
of $Y_M$ and $b_1,\ldots,b_{99},w$ restricted to $C_p$ generates
$\mathcal R_B$ and has cost $1+100p$.
\end{lemma}

\begin{proof}
In $X$, put
\[
 N_p=\{x\in X:x(t^n)\geq p\text{ for every }n\in\mathbb Z\},
 \qquad
 C_p=(\{x\in X:x(1)<p\}\cup N_p)\times\mathbb Z/M\mathbb Z.
\]
The set $N_p$ is Borel, invariant under $t$, and null: the coordinates
$x(t^n)$ are independent, and the group elements $t^n$ are distinct.
Indeed, the probability that the condition holds for $0\leq n<N$
is $(1-p)^N$.  Thus $\nu_M(C_p)=p$.
If a $t$-orbit does not meet the first part of $C_p$, all its points
belong to $N_p\times\mathbb Z/M\mathbb Z$.  Hence $C_p$ meets every
$t$-orbit, including the exceptional ones.

Let $j$ be one of $b_1,\ldots,b_{99},w$, and let $z\in Y_M$.
Choose $n\in\mathbb Z$ with $zt^n\in C_p$.  The graphing contains
the path
\[
 z\longrightarrow zt^n\longrightarrow zt^nj
   \longrightarrow zt^njt^{-n}=zj,
\]
where the first and last portions use $t$ or $t^{-1}$.  The final
identity uses $jt=tj$ in $B$.  Thus every generator of $B$ is
available everywhere in the generated relation.  Conversely, all
the graphing's maps are restrictions of elements of $B$.
Each of the $100$ restricted generators costs $p$, and the full
$t$-map costs one.
\end{proof}

\begin{proposition}
\label{actions:low}
For every integer $M>100$, the free p.m.p. action
$\Gamma\curvearrowright(Y_M,\nu_M)$ of
\eqref{actions:product-action} satisfies
\[
 \operatorname{Cost}(\mathcal R_{\Gamma\curvearrowright Y_M})
 \leq1+\frac{99}{M}.
\]
\end{proposition}

\begin{proof}
Fix $0<p<1$ and use the graphing of
Lemma~\ref{actions:small-marker}.  Add the restriction of $a$ to
the union of the height fibers $0,1,\ldots,98$.  Its domain has
measure $99/M$, so the resulting graphing has cost
\[
 1+100p+\frac{99}{M}.
\]
Let $E$ be the equivalence relation it generates.  It contains
$\mathcal R_B$, and we will prove that it contains every $a$-edge.

Suppose that $(z,za)\in E$ whenever the height of $z$ is one of
$k,k+1,\ldots,k+98$, interpreted modulo $M$.  Given $z$ at height
$k+99$, set $y=zu_{99}^{-1}$, which has height $k$.
The sequence
\[
 y=yu_0,\quad yu_1,\quad\ldots,\quad yu_{99}=z
\]
is connected by $E$: at its $i$th step apply $a$ to $yu_{i-1}$,
of height $k+i-1$, and then apply $b_i$.  The $a$-edge is among
those assumed available, and the $b_i$-edge belongs to
$\mathcal R_B$.  The $w$-edge also belongs to $\mathcal R_B$, and
\[
 yw=yu_{99}a=za.
\]
Thus $z$ and $za$ are $E$-equivalent, proving the implication from
$99$ consecutive known source heights to the next one.

The heights $0,\ldots,98$ are known initially.  Apply this
implication for the ordinary integers
$k=0,1,\ldots,M-100$.  It successively supplies source heights
$99,100,\ldots,M-1$; the last $a$-edges end at height $0$ modulo
$M$, as required.  Hence all $a$-edges lie in $E$.
Since $B$ and $a$ generate $\Gamma$, this graphing generates
$\mathcal R_\Gamma$.  Taking the infimum over graphings and then
letting $p\downarrow0$ proves the bound.
\end{proof}
